% GENERATED FILE. Edit canonical lessons, then run npm run generate:books.
% canonical-source-hash: fnv1a64:0586d038a59317a6
% canonical-lessons: MR-C418-nagar, MR-C418-mahanagar, MR-C418-nagarpalika, MR-C418-khede, MR-C418-taluka

\chapter{Town and Metropolis}
\label{ch:mr-town-and-metropolis}
\hlchaptermodality{418}

\begin{chapteropening}
\textbf{By the end of this chapter:} \emph{I can talk about a town, a big city, the town council, a hamlet and a taluka.}

Complete the last lesson of chapter 418: I can talk about a town, a big city, the town council, a hamlet and a taluka.
\end{chapteropening}

\section[\texorpdfstring{\mr{नगर} (nagar)}{नगर (nagar)}]{\mr{नगर} (nagar) — a town}

\label{lesson:MR-C418-nagar}



\begin{quote}
Before the new one: say the Marathi for a strongbox, then the Marathi for an elevator.
\end{quote}

\subsection*{You'll want to know: \mr{नगर}}
\textbf{\mr{नगर}} — \emph{nagar} — \textquotedblleft{}a town\textquotedblright{}.

\begin{quote}
\mr{हे} \mr{नगर} \mr{नदीच्या} \mr{काठावर} \mr{वसलं} \mr{आहे}.
\end{quote}

\begin{quote}
\emph{he nagar nadīcyā kāṭhāvar vaslaṁ āhe.}
\end{quote}

\begin{quote}
This town stands on the bank of the river.
\end{quote}

Neuter.

\mr{नगर} also ends many town names.

\begin{cousinweb}
\textbf{\mr{नगर}} (\emph{nagar}) is a tatsama, taken from Sanskrit \emph{nagara}, \textquotedblleft{}a town, a city\textquotedblright{} (Monier-Williams).
\end{cousinweb}

\begin{grammarlens}[title={What you've built}]
One more word for getting around.
\end{grammarlens}

\subsection*{Your turn}
\begin{itemize}
\raggedright
  \item \emph{Say it:} \emph{nagar}
  \item \emph{Say it:} \emph{nagar}, once more
  \item \emph{Say it:} \emph{he nagar nadīcyā kāṭhāvar vaslaṁ āhe.} — This town stands on the bank of the river.
  \item \emph{Recall:} say the Marathi for a strongbox, then the Marathi for an elevator, then say \emph{nagar} again
\end{itemize}

\begin{tcolorbox}[breakable,colback=teal!4,colframe=teal!35!black,title={Before you move on}]
What does \mr{नगर} mean? (\textquotedblleft{}A town\textquotedblright{}.) Say it once more.
\end{tcolorbox}
\section[\texorpdfstring{\mr{महानगर} (mahānagar)}{महानगर (mahānagar)}]{\mr{महानगर} (mahānagar) — a metropolis}

\label{lesson:MR-C418-mahanagar}



\begin{quote}
Before the new one: say the Marathi for an elevator, then the Marathi for a town.
\end{quote}

\subsection*{You'll want to know: \mr{महानगर}}
\textbf{\mr{महानगर}} — \emph{mahānagar} — \textquotedblleft{}a metropolis\textquotedblright{}.

\begin{quote}
\mr{पुणे} \mr{आता} \mr{एक} \mr{मोठं} \mr{महानगर} \mr{आहे}.
\end{quote}

\begin{quote}
\emph{puṇe ātā ek moṭhaṁ mahānagar āhe.}
\end{quote}

\begin{quote}
Pune is now a big metropolis.
\end{quote}

Neuter.

\mr{महा}-, great: a great city.

\begin{cousinweb}
\textbf{\mr{महानगर}} (\emph{mahānagar}) is a tatsama, taken from Sanskrit \emph{mahānagara}, \textquotedblleft{}a great city\textquotedblright{} (Monier-Williams).
\end{cousinweb}

\begin{grammarlens}[title={What you've built}]
One more word for getting around.
\end{grammarlens}

\subsection*{Your turn}
\begin{itemize}
\raggedright
  \item \emph{Say it:} \emph{mahānagar}
  \item \emph{Say it:} \emph{mahānagar}, once more
  \item \emph{Say it:} \emph{puṇe ātā ek moṭhaṁ mahānagar āhe.} — Pune is now a big metropolis.
  \item \emph{Recall:} say the Marathi for an elevator, then the Marathi for a town, then say \emph{mahānagar} again
\end{itemize}

\begin{tcolorbox}[breakable,colback=teal!4,colframe=teal!35!black,title={Before you move on}]
What does \mr{महानगर} mean? (\textquotedblleft{}A metropolis\textquotedblright{}.) Say it once more.
\end{tcolorbox}
\section[\texorpdfstring{\mr{नगरपालिका} (nagarpālikā)}{नगरपालिका (nagarpālikā)}]{\mr{नगरपालिका} (nagarpālikā) — a municipal council}

\label{lesson:MR-C418-nagarpalika}



\begin{quote}
Before the new one: say the Marathi for a town, then the Marathi for a metropolis.
\end{quote}

\subsection*{You'll want to know: \mr{नगरपालिका}}
\textbf{\mr{नगरपालिका}} — \emph{nagarpālikā} — \textquotedblleft{}a municipal council\textquotedblright{}.

\begin{quote}
\mr{नगरपालिकेने} \mr{नवीन} \mr{रस्ता} \mr{बांधला}.
\end{quote}

\begin{quote}
\emph{nagarpālikene navīn rastā bāndhlā.}
\end{quote}

\begin{quote}
The municipal council built a new road.
\end{quote}

Feminine.

Before an ending -\mr{का} becomes -\mr{के}-: \mr{नगरपालिकेने}.

\begin{cousinweb}
\textbf{\mr{नगरपालिका}} (\emph{nagarpālikā}) is a modern word built from Sanskrit \emph{nagara}, \textquotedblleft{}a town\textquotedblright{}, and \emph{pālaka}, \textquotedblleft{}a guardian, a protector\textquotedblright{} (Monier-Williams), here in its feminine \emph{pālikā}.
\end{cousinweb}

\begin{grammarlens}[title={What you've built}]
One more word for getting around.
\end{grammarlens}

\subsection*{Your turn}
\begin{itemize}
\raggedright
  \item \emph{Say it:} \emph{nagarpālikā}
  \item \emph{Say it:} \emph{nagarpālikā}, once more
  \item \emph{Say it:} \emph{nagarpālikene navīn rastā bāndhlā.} — The municipal council built a new road.
  \item \emph{Recall:} say the Marathi for a town, then the Marathi for a metropolis, then say \emph{nagarpālikā} again
\end{itemize}

\begin{tcolorbox}[breakable,colback=teal!4,colframe=teal!35!black,title={Before you move on}]
What does \mr{नगरपालिका} mean? (\textquotedblleft{}A municipal council\textquotedblright{}.) Say it once more.
\end{tcolorbox}
\section[\texorpdfstring{\mr{खेडे} (kheḍe)}{खेडे (kheḍe)}]{\mr{खेडे} (kheḍe) — a small village; a hamlet}

\label{lesson:MR-C418-khede}



\begin{quote}
Before the new one: say the Marathi for a metropolis, then the Marathi for a municipal council.
\end{quote}

\subsection*{You'll want to know: \mr{खेडे}}
\textbf{\mr{खेडे}} — \emph{kheḍe} — \textquotedblleft{}a small village; a hamlet\textquotedblright{}.

\begin{quote}
\mr{माझे} \mr{आजोबा} \mr{एका} \mr{लहान} \mr{खेड्यात} \mr{राहतात}.
\end{quote}

\begin{quote}
\emph{mājhe ājobā ekā lahān kheḍyāt rāhtāt.}
\end{quote}

\begin{quote}
My grandfather lives in a small village.
\end{quote}

Neuter. Plural: \mr{खेडी}, \emph{kheḍī}.

A neuter noun: before an ending \mr{खेड्या}-, in the plural \mr{खेडी}.

\begin{cousinweb}
\textbf{\mr{खेडे}} (\emph{kheḍe}) is inherited: CDIAL lists Marathi \emph{kheḍẽ}, \textquotedblleft{}hamlet\textquotedblright{}, under Sanskrit \emph{khēṭa}, \textquotedblleft{}village\textquotedblright{} (CDIAL 3916).
\end{cousinweb}

\begin{grammarlens}[title={What you've built}]
One more word for getting around.
\end{grammarlens}

\subsection*{Your turn}
\begin{itemize}
\raggedright
  \item \emph{Say it:} \emph{kheḍe}
  \item \emph{Say it:} \emph{kheḍe}, once more
  \item \emph{Say it:} \emph{mājhe ājobā ekā lahān kheḍyāt rāhtāt.} — My grandfather lives in a small village.
  \item \emph{Recall:} say the Marathi for a metropolis, then the Marathi for a municipal council, then say \emph{kheḍe} again
\end{itemize}

\begin{tcolorbox}[breakable,colback=teal!4,colframe=teal!35!black,title={Before you move on}]
What does \mr{खेडे} mean? (\textquotedblleft{}A small village; a hamlet\textquotedblright{}.) Say it once more.
\end{tcolorbox}
\section[\texorpdfstring{\mr{तालुका} (tālukā)}{तालुका (tālukā)}]{\mr{तालुका} (tālukā) — a taluka}

\label{lesson:MR-C418-taluka}



\begin{quote}
Before the new one: say the Marathi for a municipal council, then the Marathi for a small village.
\end{quote}

\subsection*{You'll want to know: \mr{तालुका}}
\textbf{\mr{तालुका}} — \emph{tālukā} — \textquotedblleft{}a taluka\textquotedblright{}.

\begin{quote}
\mr{आमचा} \mr{तालुका} \mr{डोंगराळ} \mr{आहे}.
\end{quote}

\begin{quote}
\emph{āmcā tālukā ḍoṅgrāḷ āhe.}
\end{quote}

\begin{quote}
Our taluka is hilly.
\end{quote}

Masculine.

A part of a district, with its own town office; the plural is the same.

\begin{cousinweb}
\textbf{\mr{तालुका}} (\emph{tālukā}) is a loan through Persian, usually traced to Arabic \emph{taʿalluqa}, \textquotedblleft{}a dependency, an estate\textquotedblright{}.
\end{cousinweb}

\begin{grammarlens}[title={What you've built}]
One more word for getting around.
\end{grammarlens}

\subsection*{Your turn}
\begin{itemize}
\raggedright
  \item \emph{Say it:} \emph{tālukā}
  \item \emph{Say it:} \emph{tālukā}, once more
  \item \emph{Say it:} \emph{āmcā tālukā ḍoṅgrāḷ āhe.} — Our taluka is hilly.
  \item \emph{Recall:} say the Marathi for a municipal council, then the Marathi for a small village, then say \emph{tālukā} again
\end{itemize}

\begin{tcolorbox}[breakable,colback=teal!4,colframe=teal!35!black,title={Before you move on}]
What does \mr{तालुका} mean? (\textquotedblleft{}A taluka\textquotedblright{}.) Say it once more.
\end{tcolorbox}
